% TOP_PROBLEM: {"id": 5500049, "title": "Planar Euclidean Maximum TSP", "category_id": 6, "category": {"id": 6, "name": "geometry", "display_name": "Geometry", "description": "Euclidean and non-Euclidean geometry, geometric structures.", "slug": "geometry", "order_index": 6, "created_at": "2026-07-31T15:26:25.671Z"}, "status": "open", "problem_number": "AMR-054-0049", "set_id": 13}
% FIRST_POSTED: 2026-10-02
% ENTRY_KIND: statement_only
% UnsolvedMath ID 5500049; source catalogue code AMR-054-0049
% !TeX program = lualatex
% Definition and sources only; this file is not an LLM solution attempt.
\documentclass[11pt,a4paper]{article}
\usepackage[margin=25mm]{geometry}
\usepackage{fontspec}
\setmainfont{lmroman10-regular.otf}[BoldFont=lmroman10-bold.otf,
 ItalicFont=lmroman10-italic.otf,BoldItalicFont=lmroman10-bolditalic.otf]
\usepackage{amsmath,amssymb,mathtools}
\usepackage{unicode-math}
\setmathfont{latinmodern-math.otf}
\AtBeginDocument{\let\setminus\smallsetminus}
\usepackage{xurl,hyperref}
\hypersetup{colorlinks=true,urlcolor=blue}
\setcounter{secnumdepth}{0}
\setlength{\parindent}{0pt}
\setlength{\parskip}{5pt}
\setlength{\emergencystretch}{3em}
\begin{document}
\section*{Planar Euclidean Maximum TSP}\label{5500049}
{\small UnsolvedMath ID 5500049 · AMR-054-0049 · Geometry · AMR Open Problem Lists}
\subsection{Definitions and mathematical statement}
Given $n\ge3$ distinct points $P=\{p_1,\ldots,p_n\}$ in
$\mathbb R^2$, a tour is a cyclic permutation $\sigma$ of all
points, with Euclidean length
\[
  \mathcal L(\sigma)=\sum_{i=1}^{n}
          \|p_{\sigma(i)}-p_{\sigma(i+1)}\|_2,
                    \qquad\sigma(n+1)=\sigma(1).
\]
The planar Euclidean maximum traveling-salesperson problem
asks for a tour maximizing this sum. What is the computational
complexity of this exact optimization problem: does it have
a polynomial-time algorithm, or can an appropriate hardness
result be established?

The tour may cross itself; imposing simplicity would change
the feasible set. All points are visited once before returning
to the start, and the objective is total length, not the longest
single edge or a bottleneck statistic. The dimension is fixed
at two, and distances come from the Euclidean norm.

A complexity statement must specify its numerical model.
One may use exact real arithmetic for geometric operations;
with rational coordinates in a bit model, comparisons of sums
of square roots also need to be accounted for. The question
concerns exact optimality. A polynomial-time approximation
scheme, or an algorithm for maximum TSP in a different metric
or dimension, does not by itself determine this planar exact
complexity.
\subsection{Short English statement}
Determine the complexity of finding a maximum-length Euclidean Hamiltonian tour through a planar point set.
\subsection{Sources}
\begin{itemize}
\item[S1] ULAM.AI and the UnsolvedMath Contributors, supplied problems.json, record 5500049 (AMR-054-0049), AMR Open Problem Lists. Catalogue identity and source transcription; CC BY 4.0.
\url{https://huggingface.co/datasets/ulamai/UnsolvedMath}.
\item[S2] Open Problems Project - Computational Geometry, Problem 49. Original source indexed by AMR-054-0049.
\url{https://topp.openproblem.net/p49}.
\item[S3] The Open Problems Project, original numbered problem and status notes.
\url{https://topp.openproblem.net/p49}.
\end{itemize}
\end{document}
